% ==============================================================================
% FFGFT: Document 245 (EN)
% Title: FFGFT in the RA 2.1 Architecture — Diagnostic Entry
% Author: Johann Pascher · May 2026
% RA nomenclature: Guevara Calderón, J.T. (2026). RA 2.1. IPI working document.
% ==============================================================================

\subsection*{List of Symbols}

{\small
\begin{tabular}{p{3.2cm}p{10.5cm}}
\toprule
\textbf{Symbol} & \textbf{Meaning} \\
\midrule
$\xi = 4/30000$ & FFGFT dimensionless winding parameter; geometrically fixed via $\kappa=7$ (not fitted) \\
$T^4$ & 4-dimensional torus $(\mathbb{R}/2\pi\mathbb{Z})^4$; FFGFT ontological primitive \\
$n_\theta, n_\phi, n_3, n_4$ & Integer winding numbers in the four torus directions \\
$\theta_4$ & Continuous fourth-phase coordinate on $T^4$ \\
$\varepsilon \approx 0.001$ & Non-closure of $T^4$; generative structural residual \\
$\tilde{T}\cdot m = 1$ & Foundational T0 relation (time-mass duality) \\
$D_f = 3 - \xi$ & Fractal dimension of physical space ($\approx 2.99987$) \\
$\kappa = 7$ & Unique winding exponent from three measured mass ratios \\
$L^2(T^4)$ & FFGFT Hilbert space; square-integrable functions on $T^4$ \\
$\alpha$ & Fine structure constant; derived from $\xi$ (Doc~009/137) \\
$\Lambda_\text{FFGFT}$ & Cosmological constant: $1.34\times10^{-122}$ Planck units (Doc~182) \\
$H_0$ & Hubble constant from $\xi$: $66.2\,\text{km/s/Mpc}$, $-1.9\%$ from Planck (Doc~182) \\
$\tau\cdot\mu$ & Coherence-time--information-mass product $= \hbar/c^2$; exact \\
RA4--RA1 & Layers of the Representational Architecture, introduced by J.T.~Guevara Calderón (2026), \textit{RA 2.1}, IPI working document, May 2026 \\
\bottomrule
\end{tabular}
}

\vspace{0.6em}\hrule\vspace{0.5em}

\begin{center}
{\large\bfseries FFGFT Diagnostic Entry — RA 2.1 Format}\\[0.2em]
{\normalsize Fundamental Fractal Geometric Field Theory (T0 / Time-Mass Duality)}\\[0.1em]
{\small Johann Pascher · ORCID 0009-0000-6518-4064 · May 2026}\\[0.1em]
{\footnotesize RA nomenclature: Guevara Calderón, J.T. (2026). \textit{Representational Architecture 2.1}. IPI working document, May 2026.}
\end{center}

\vspace{0.3em}\hrule\vspace{0.3em}

\renewcommand{\arraystretch}{1.4}
{\small
\begin{longtable}{p{2.0cm}p{4.4cm}p{4.4cm}p{2.55cm}}

\toprule
\textbf{Layer} & \textbf{FFGFT Content} & \textbf{Standard Model / \(\Lambda\)CDM} & \textbf{RA Status} \\
\midrule
\endfirsthead
\toprule
\textbf{Layer} & \textbf{FFGFT Content} & \textbf{Standard Model / \(\Lambda\)CDM} & \textbf{RA Status} \\
\midrule
\endhead
\midrule\multicolumn{4}{r}{\footnotesize\itshape continued\ldots}\\\endfoot
\bottomrule\endlastfoot
\sloppy\raggedright

\textbf{RA4}\\Ontol. &
4D torus $T^4$ with winding structure $(n_\theta, n_\phi, n_3, n_4)$, continuous phase $\theta_4$, and non-closure $\varepsilon \approx 0.001$. Single parameter $\xi = 4/30000$ geometrically fixed from three measured mass ratios via $\kappa = 7$. Foundational relation $\tilde{T}\cdot m = 1$. &
Quantum fields on Minkowski spacetime $\mathbb{R}^{3,1}$. Multiple independent coupling constants, masses, and mixing angles as inputs. $\Lambda$ identified with QFT vacuum energy (Zeldovich step~3). &
\textbf{RA4 divergence.} Primitive replaced. Zeldovich step~3 not made; no tuning problem. \\

\midrule

\textbf{RA3}\\Admiss. &
Only winding configurations consistent with $\xi$ are physically admissible. Fractal dimension $D_f = 3-\xi \approx 2.99987$ defines accessible topological sectors. Z3$^3$ topological constraint restricts to exactly three particle generations. No additional axioms needed. &
All renormalizable field configurations admissible. Particle content and number of generations are experimental inputs, not derived from any admissibility constraint. &
\textbf{RA3 restriction.} $\xi$ maximally constrains admissibility. SM leaves it open; generations unexplained. \\

\midrule

\textbf{RA2}\\Repres. &
$L^2(T^4)$ as state space. Fourier modes with winding numbers \mbox{$\mathbf{n}\in\mathbb{Z}^4$} as natural basis. Standard QM operator algebra and Born rule emerge as projections from torus geometry (Doc~230/231,
  Hilbert space bridge). No separate quantization postulate. &
Fock space over Minkowski spacetime. Operator algebra and Born rule postulated independently of spacetime structure. Quantization is an additional input. &
\textbf{RA2 derivation.} QM formalism derived from $T^4$ geometry. In SM it is postulated independently. \\

\midrule

\textbf{RA1.5}\\Operat. &
All derived from $\xi$ alone, \textbf{no fitting}:
{\small\begin{itemize}[leftmargin=*,nosep,itemsep=1pt]
  \item Lepton and quark masses via $\kappa=7$ recursion; agreement $<1\%$
  \item Fine structure constant $\alpha$ (Doc~009/137)
  \item CMB temperature; deviation $\Delta = 0.0002\%$ (Doc~025)
  \item $\Lambda_\text{FFGFT} = 1.34\times10^{-122}$ Planck units (Doc~182)
  \item $H_0 = 66.2\,\text{km/s/Mpc}$ ($-1.9\%$ from Planck)
  \item $\tau\cdot\mu = \hbar/c^2$ (exact)
\end{itemize}} &
All quantities above are independent experimental inputs requiring separate measurement. $\Lambda$ requires 120-order fine-tuning relative to QFT zero-point energy. &
\textbf{RA1.5 unification.} All SM constants from one parameter. SM treats them as unrelated experimental inputs. \\

\midrule

\textbf{RA1}\\Expr. &
Concrete numerical values for all Standard Model constants and cosmological parameters, derived not fitted. Published in 149+ documents; GitHub: jpascher/T0-Time-Mass-Duality; Zenodo v1.1.0, DOI 10.5281/zenodo.20117635. &
Parameter values measured and inserted by hand. No unifying derivation. SM has $\approx\!20$ free parameters; $\Lambda$CDM adds several more. &
\textbf{RA1 compression.} 1 parameter vs.\ $\sim$20 (SM). \\

\end{longtable}
}

\vspace{0.3em}\hrule\vspace{0.4em}

\noindent\textbf{Falsification criteria (RA1 level):}
\begin{itemize}[nosep,itemsep=1pt]
  \item Any lepton or quark mass ratio inconsistent with the $\xi$-ladder at $\kappa=7$ (tolerance $<1\%$).
  \item CMB deviation exceeding the correction term $(275/4)\cdot\xi \approx 9.2\times10^{-3}$ (Doc~025).
  \item A cosmological constant value outside the range derived in Doc~182.
  \item Any winding quantum not assignable to an integer in $\mathbb{Z}^4$.
\end{itemize}

\vspace{0.3em}
\noindent\textbf{Key references:}
Doc~009 ($\xi$, $\kappa=7$) ·
Doc~013 (SI units) ·
Doc~025 (CMB) ·
Doc~182 ($\Lambda$, $H_0$) ·
Doc~230/231 (Hilbert space) ·
\href{https://github.com/jpascher/T0-Time-Mass-Duality}{GitHub} ·
\href{https://doi.org/10.5281/zenodo.20117635}{Zenodo v1.1.0}
